% =============================================================
%  MechanicPro --- User Manual
%  Template: Professional / Modern / Clean
% =============================================================

\documentclass[12pt, a4paper]{report}

% -- Encoding & Font ------------------------------------------
\usepackage[T1]{fontenc}
\usepackage[utf8]{inputenc}
\usepackage[scaled=0.92]{helvet}
\renewcommand{\familydefault}{\sfdefault}

% -- Page Layout ----------------------------------------------
\usepackage[
top=2.5cm, bottom=2.5cm,
left=3cm,  right=2.5cm
]{geometry}
% Automatic "Month Year" date
\newcommand{\monthyear}{%
	\ifcase\month\or January\or February\or March\or April\or May\or June\or
	July\or August\or September\or October\or November\or December\fi
	\space\the\year%
}

% -- Colors ---------------------------------------------------
\usepackage{xcolor}
\definecolor{Navy}{RGB}{15, 40, 75}
\definecolor{Blue}{RGB}{37, 99, 195}
\definecolor{LightBlue}{RGB}{219, 234, 254}
\definecolor{TextGray}{RGB}{90, 105, 120}
\definecolor{LightGray}{RGB}{248, 249, 251}
\definecolor{RuleGray}{RGB}{220, 225, 232}
\definecolor{StatusGreen}{RGB}{22, 163, 74}
\definecolor{WarnYellow}{RGB}{254, 243, 199}
\definecolor{WarnBorder}{RGB}{217, 119, 6}
\definecolor{DangerRed}{RGB}{254, 226, 226}
\definecolor{DangerBorder}{RGB}{185, 28, 28}

% -- Graphics & Drawing ---------------------------------------
\usepackage{graphicx}
\usepackage{tikz}
\usetikzlibrary{calc}

% -- Tables ---------------------------------------------------
\usepackage{booktabs}
\usepackage{tabularx}
\usepackage{array}
\usepackage{longtable}

% -- Lists ----------------------------------------------------
\usepackage{enumitem}
\setlist[itemize]{leftmargin=1.5em, itemsep=3pt}
\setlist[enumerate]{leftmargin=1.5em, itemsep=3pt}

% Compact step list (numbered, used for procedures)
\newlist{steps}{enumerate}{1}
\setlist[steps]{
	label=\textbf{\color{Blue}\arabic*.},
	leftmargin=2em,
	itemsep=5pt,
	topsep=6pt,
}

% -- Headers & Footers ----------------------------------------
\usepackage{fancyhdr}
\pagestyle{fancy}
\fancyhf{}
\fancyhead[L]{%
	\small\color{TextGray}\textit{MechanicPro \textemdash{} User Manual}%
}
\fancyhead[R]{%
	\small\color{TextGray}v1.2.0%
}
\fancyfoot[L]{%
	\small\color{TextGray}User Manual%
}
\fancyfoot[R]{%
	\small\color{TextGray}\thepage%
}
\renewcommand{\headrulewidth}{0.4pt}
\renewcommand{\footrulewidth}{0.4pt}
\renewcommand{\headrule}{%
	\color{RuleGray}\hrule width\headwidth height\headrulewidth%
}
\renewcommand{\footrule}{%
	\color{RuleGray}\hrule width\headwidth height\footrulewidth%
}

% Plain style for chapter-opening pages
\fancypagestyle{plain}{%
	\fancyhf{}
	\fancyfoot[R]{\small\color{TextGray}\thepage}
	\renewcommand{\headrulewidth}{0pt}
	\renewcommand{\footrulewidth}{0.4pt}
	\renewcommand{\footrule}{%
		\color{RuleGray}\hrule width\headwidth height\footrulewidth%
	}
}

% -- Section Formatting ---------------------------------------
\usepackage{titlesec}

\titleformat{\chapter}[block]
{\normalfont\huge\bfseries\color{Navy}}
{\color{Blue}\thechapter}{1em}{}
[\vspace{4pt}\color{Blue}\rule{\textwidth}{1.5pt}\vspace{-6pt}]

\titlespacing{\chapter}{0pt}{0pt}{20pt}

\titleformat{\section}
{\normalfont\Large\bfseries\color{Navy}}
{\color{Blue}\thesection}{1em}{}

\titleformat{\subsection}
{\normalfont\large\bfseries\color{TextGray}}
{\thesubsection}{1em}{}

\titleformat{\subsubsection}
{\normalfont\normalsize\bfseries\color{TextGray}}
{\thesubsubsection}{1em}{}

% -- TOC Styling ----------------------------------------------
\usepackage{tocloft}
\renewcommand{\cftchapfont}{\bfseries\color{Navy}}
\renewcommand{\cftchappagefont}{\bfseries\color{Navy}}
\renewcommand{\cftsecfont}{\color{TextGray}}
\renewcommand{\cftsecpagefont}{\color{TextGray}}
\setlength{\cftbeforechapskip}{6pt}

% -- Hyperlinks -----------------------------------------------
\usepackage{hyperref}
\hypersetup{
	colorlinks  = true,
	linkcolor   = Blue,
	urlcolor    = Blue,
	citecolor   = Blue,
	pdftitle    = {MechanicPro --- User Manual},
	pdfauthor   = {Your Name},
	pdfsubject  = {User Manual},
	pdfkeywords = {MechanicPro, mechanic shop, user manual},
}

% -- Custom Commands ------------------------------------------

% Note box (blue)
\newcommand{\notebox}[1]{%
	\vspace{6pt}
	\noindent
	\colorbox{LightBlue}{%
		\begin{minipage}{\dimexpr\linewidth-2\fboxsep}%
			\vspace{5pt}
			\small\textbf{\color{Blue}Note:\space}\color{Navy}#1
			\vspace{5pt}
		\end{minipage}%
	}
	\vspace{6pt}
}

% Warning box (yellow)
\newcommand{\warnbox}[1]{%
	\vspace{6pt}
	\noindent
	\colorbox{WarnYellow}{%
		\begin{minipage}{\dimexpr\linewidth-2\fboxsep}%
			\vspace{5pt}
			\small\textbf{\color{WarnBorder}Warning:\space}\color{Navy}#1
			\vspace{5pt}
		\end{minipage}%
	}
	\vspace{6pt}
}

% Danger box (red) -- for destructive actions
\newcommand{\dangerbox}[1]{%
	\vspace{6pt}
	\noindent
	\colorbox{DangerRed}{%
		\begin{minipage}{\dimexpr\linewidth-2\fboxsep}%
			\vspace{5pt}
			\small\textbf{\color{DangerBorder}Caution:\space}\color{Navy}#1
			\vspace{5pt}
		\end{minipage}%
	}
	\vspace{6pt}
}

% UI element label (buttons, menus, field names)
\newcommand{\ui}[1]{\textbf{\color{Navy}#1}}

% Keyboard shortcut
\newcommand{\key}[1]{%
	\tikz[baseline=(k.base)]\node[
	draw=RuleGray, fill=LightGray, rounded corners=2pt,
	inner xsep=4pt, inner ysep=2pt, font=\small\ttfamily
	](k){#1};%
}

% Placeholder for screenshots
\newcommand{\screenshot}[1]{%
	\begin{center}
		\begin{tikzpicture}
			\draw[dashed, color=RuleGray, fill=LightGray, rounded corners=4pt]
			(0,0) rectangle (\linewidth-0.4cm, 5cm);
			\node[color=TextGray] at (0.5\linewidth-0.2cm, 2.5cm)
			{\small\textit{[Screenshot: #1]}};
		\end{tikzpicture}
	\end{center}
}

% -- Abstract Page --------------------------------------------
\newcommand{\abstractpage}[1]{%
	\clearpage
	\thispagestyle{plain}
	\vspace*{2cm}
	\begin{center}
		{\color{Navy}\Large\bfseries Abstract}\\[0.4cm]
		\color{Blue}\rule{0.3\textwidth}{1.5pt}
	\end{center}
	\vspace{0.8cm}
	\begin{quote}
		\setlength{\parskip}{6pt}
		\normalsize #1
	\end{quote}
	\vfill
	\noindent\small\color{TextGray}
	\textbf{Keywords:} MechanicPro, user manual, mechanic shop,
	service orders, clients, vehicles, CSV, PDF.
	\clearpage
}

% -- Cover Page -----------------------------------------------
\newcommand{\coverpage}{%
	\begin{titlepage}
		\begin{tikzpicture}[remember picture, overlay]
			
			%% Top bar
			\fill[Navy]
			(current page.north west)
			rectangle
			([yshift=-4cm] current page.north east);
			
			%% Accent stripe
			\fill[Blue]
			([yshift=-4cm] current page.north west)
			rectangle
			([yshift=-4.12cm] current page.north east);
			
			%% Bottom band
			\fill[LightGray]
			(current page.south west)
			rectangle
			([yshift=3.2cm] current page.south east);
			
			%% Bottom band top border
			\fill[RuleGray]
			([yshift=3.2cm] current page.south west)
			rectangle
			([yshift=3.24cm] current page.south east);
			
		\end{tikzpicture}
		
		%% -- Top bar content -----------------------
		\vspace*{1.2cm}
		\hspace{1.2cm}
		\begin{minipage}{0.8\textwidth}
			{\color{white}\footnotesize\bfseries\MakeUppercase{User Manual}}
			\quad
			{\color{Blue!40!white}\footnotesize User Manual}
		\end{minipage}
		
		\vspace{3.6cm}
		
		%% -- Main title block ----------------------
		\hspace{1.2cm}
		\begin{minipage}{0.82\textwidth}
			
			{\color{Navy}\Huge\bfseries\scalebox{2.0}{MechanicPro}}
			
			\vspace{0.3cm}
			{\color{TextGray}\large Management System for a Mechanic Shop}
			
			\vspace{0.6cm}
			\color{Blue}\rule{\textwidth}{1.5pt}
			\vspace{0.3cm}
			
			{\color{TextGray}\normalsize
				Complete Guide to Installation, Daily Use, and Data Management
			}
			
		\end{minipage}
		
		\vfill
		
		%% -- Metadata block ------------------------
		\hspace{1.2cm}
		\begin{minipage}{0.82\textwidth}
			\renewcommand{\arraystretch}{1.5}
			\begin{tabular}{>{\color{TextGray}\small\bfseries}l
					>{\small}l
					@{\hskip 2cm}
					>{\color{TextGray}\small\bfseries}l
					>{\small}l}
				Author    & Your Name      & Version   & 1.2.0      \\
				Date      & \monthyear     & Status    &
				\textcolor{StatusGreen}{\bfseries Final}        \\
				Document  & User Manual & Language  & English    \\
			\end{tabular}
		\end{minipage}
		
		\vspace{2.8cm}
		
	\end{titlepage}
}

% =============================================================
%  DOCUMENT
% =============================================================

\begin{document}
	
	\coverpage
	
	\abstractpage{%
		This document is the user manual for \textbf{MechanicPro}, a desktop
		application for managing the operations of a mechanic shop. It covers
		everything a user needs to install and run the application, navigate its
		interface, and perform all available tasks.
		
		The manual is organised by workflow. Each chapter corresponds to a
		distinct area of the application: managing clients and vehicles, handling
		service orders, using the dashboard, and working with data import,
		export, and PDF generation.
		
		No technical knowledge is required to follow this manual. For information
		about the system architecture or database design, refer to the Technical
		Report (TR-MECHPRO-001).
	}
	
	% -- Revision History -----------------------------------------
	\chapter*{Revision History}

	\noindent \begin{tabularx}{\textwidth}{c l X l}
		\toprule
		\textbf{Version} & \textbf{Date} & \textbf{Description} & \textbf{Author} \\
		\midrule
		1.0.0 & Jan 2026 & Initial document created          & Your Name \\
		1.2.0 & May 2026 & CSV and PDF sections added        & Your Name \\
		\bottomrule
	\end{tabularx}
	
	\newpage
	
	\tableofcontents
	
	\newpage
	
	\listoffigures
	
	\newpage
	
	\listoftables
	
	
	% =============================================================
	%  CHAPTER 1 --- INTRODUCTION
	% =============================================================
	\chapter{Introduction}
	
	\section{About This Manual}
	% Explain the purpose of this document and how it is structured.
	% Who it is for, what it covers, and what it does not cover.
	
	\section{About MechanicPro}
	% Brief, non-technical description of what the application does.
	% Focus on the user's perspective: what problems it solves.
	
	\section{Conventions Used in This Manual}
	
	Throughout this document the following conventions are used:
	
	\begin{itemize}
		\item \ui{Bold navy text} refers to interface elements such as buttons,
		menu items, and field labels.
		\item \key{Ctrl+S} represents keyboard shortcuts.
		\item Steps to complete a task are numbered and shown like this:
	\end{itemize}
	
	\begin{steps}
		\item This is how a step looks.
		\item Each step describes one action.
	\end{steps}
	
	\notebox{Blue boxes contain tips or additional context.}
	
	\warnbox{Yellow boxes highlight situations that may produce unexpected results.}
	
	\dangerbox{Red boxes warn about irreversible actions, such as deleting records.}
	
	% =============================================================
	%  CHAPTER 2 --- INSTALLATION
	% =============================================================
	\chapter{Installation}
	
	\section{System Requirements}
	
	\begin{tabularx}{\textwidth}{l X}
		\toprule
		\textbf{Component} & \textbf{Requirement} \\
		\midrule
		Operating System & Windows 10 or later / macOS 11 or later / Ubuntu 20.04 or later \\
		RAM              & 4 GB minimum, 8 GB recommended \\
		Disk Space       & % Fill in after checking the built binary size
		\\
		Additional       & Windows: WebView2 runtime (included in Windows 11) \\
		\bottomrule
	\end{tabularx}
	
	\section{Installing the Application}
	
	\subsection{Windows}
	
	\begin{steps}
		\item Download the \texttt{MechanicPro\_1.2.0\_x64.msi} installer.
		\item Double-click the installer file.
		\item Follow the on-screen instructions.
		\item Launch \ui{MechanicPro} from the Start Menu or the desktop shortcut.
	\end{steps}
	
	\notebox{Windows may show a security warning on first launch because the
		application is not signed with a commercial certificate. Click
		\ui{More info} and then \ui{Run anyway} to proceed.}
	
	\subsection{macOS}
	
	\begin{steps}
		\item Download the \texttt{MechanicPro\_1.2.0\_x64.dmg} file.
		\item Open the \texttt{.dmg} and drag \ui{MechanicPro} to your
		\ui{Applications} folder.
		\item On first launch, right-click the app and select \ui{Open}.
	\end{steps}
	
	\subsection{Linux}
	
	\begin{steps}
		\item Download the \texttt{.deb} or \texttt{.AppImage} package.
		\item Install with \texttt{sudo dpkg -i MechanicPro\_1.2.0\_amd64.deb}
		or make the AppImage executable and run it directly.
	\end{steps}
	
	\section{First Launch}
	% Describe what the user sees the very first time they open the app.
	% Empty state, any welcome message, where data is stored.
	
	\notebox{MechanicPro stores all data locally on your machine. No account or
		internet connection is required.}
	
	% =============================================================
	%  CHAPTER 3 --- INTERFACE OVERVIEW
	% =============================================================
	\chapter{Interface Overview}
	
	\section{Main Window}
	% Describe the overall layout: sidebar/navigation, main content area,
	% header bar, any status indicators.
	
	\screenshot{Main application window}
	
	\section{Navigation}
	% How the user moves between sections (sidebar links, top menu, etc.)
	% List all main sections accessible from navigation.
	
	\section{Dashboard}
	% Brief intro here; the full dashboard chapter comes later.
	% Just orient the user to what they see on startup.
	
	\screenshot{Dashboard overview}
	
	% =============================================================
	%  CHAPTER 4 --- CLIENTS
	% =============================================================
	\chapter{Managing Clients}
	
	\section{Viewing the Client List}
	% How to access the clients section, what the list shows,
	% how to search or filter.
	
	\screenshot{Client list view}
	
	\section{Adding a Client}
	
	\begin{steps}
		\item Navigate to the \ui{Clients} section.
		\item Click \ui{New Client}.
		\item Fill in the required fields:
		\begin{itemize}
			\item \ui{Name} --- full name of the client (required).
			\item \ui{Phone} --- contact phone number.
			\item \ui{Email} --- contact email address.
		\end{itemize}
		\item Click \ui{Save}.
	\end{steps}
	
	\section{Editing a Client}
	
	\begin{steps}
		\item Locate the client in the list.
		\item Click on the client's name or the edit icon.
		\item Update the fields as needed.
		\item Click \ui{Save}.
	\end{steps}
	
	\section{Viewing a Client's History}
	% How to see all service orders associated with a given client.
	
	% =============================================================
	%  CHAPTER 5 --- VEHICLES
	% =============================================================
	\chapter{Managing Vehicles}
	
	\section{Viewing the Vehicle List}
	% How vehicles are listed, how they relate to clients.
	
	\section{Adding a Vehicle}
	
	\begin{steps}
		\item Navigate to the \ui{Vehicles} section.
		\item Click \ui{New Vehicle}.
		\item Fill in the required fields:
		\begin{itemize}
			\item \ui{Licence Plate} --- unique identifier for the vehicle (required).
			\item \ui{Make} --- manufacturer name.
			\item \ui{Model} --- vehicle model.
			\item \ui{Year} --- year of manufacture.
			\item \ui{Owner} --- select the associated client.
		\end{itemize}
		\item Click \ui{Save}.
	\end{steps}
	
	\section{Editing a Vehicle}
	
	\begin{steps}
		\item Locate the vehicle in the list.
		\item Click the edit icon.
		\item Update the fields as needed.
		\item Click \ui{Save}.
	\end{steps}
	
	% =============================================================
	%  CHAPTER 6 --- SERVICE ORDERS
	% =============================================================
	\chapter{Service Orders}
	
	\section{Overview}
	% Explain what a service order is in the context of this app.
	% What information it holds, what states it can be in.
	
	\begin{tabularx}{\textwidth}{l X}
		\toprule
		\textbf{Status} & \textbf{Meaning} \\
		\midrule
		% Fill in the actual statuses used in the app
		Open       & The service has been registered but work has not started. \\
		In Progress & Work is currently underway. \\
		Completed  & The service has been finished. \\
		\bottomrule
	\end{tabularx}
	
	\section{Creating a Service Order}
	
	\begin{steps}
		\item Navigate to \ui{Service Orders}.
		\item Click \ui{New Order}.
		\item Select the \ui{Client} and associated \ui{Vehicle}.
		\item Describe the service requested in the \ui{Description} field.
		\item Set the \ui{Status} and any other relevant fields.
		\item Click \ui{Save}.
	\end{steps}
	
	\notebox{You can also create a new client or vehicle directly from this
		screen without leaving the service order form.}
	
	\section{Updating a Service Order}
	% How to change the status, add notes, update cost, etc.
	
	\section{Changing the Order of Services}
	% The app supports reordering --- explain how and why.
	
	\section{Viewing Service History}
	% How to access the history of completed or past orders.
	% What information is visible in history view.
	
	\screenshot{Service history view}
	
	\section{Downloading Service Details}
	% The button that generates a PDF of a service order.
	% Covered in more detail in Chapter 8.
	
	% =============================================================
	%  CHAPTER 7 --- DASHBOARD
	% =============================================================
	\chapter{Dashboard}
	
	\section{Overview}
	% What the dashboard shows at a glance.
	% Purpose: quick status check, not detailed management.
	
	\screenshot{Dashboard with data}
	
	\section{Statistics}
	% Describe each statistic shown:
	% total clients, total vehicles, open orders, completed orders, etc.
	% What the numbers represent and how they are calculated.
	
	\section{Quick Actions}
	% The shortcut buttons available on the dashboard
	% (e.g., create client, create vehicle directly from dashboard).
	
	\section{Charts}
	% Describe the charts shown (recharts library).
	% What period they cover, what they represent.
	
	% =============================================================
	%  CHAPTER 8 --- DATA MANAGEMENT
	% =============================================================
	\chapter{Data Management}
	
	\section{Exporting Data to CSV}
	% Which entities can be exported (clients, vehicles, orders).
	% How to trigger the export, where the file is saved.
	
	\begin{steps}
		\item Navigate to the \ui{Export} section.
		\item Select the data to export.
		\item Click \ui{Export CSV}.
		\item Choose a location to save the file.
	\end{steps}
	
	\section{Importing Data from CSV}
	% How to prepare a CSV file for import.
	% Column format requirements.
	% How to trigger the import.
	
	\warnbox{Importing a CSV file will add records to the existing data.
		It does not replace or delete current records.}
	
	\subsection{Required CSV Format}
	
	The import file must use UTF-8 encoding and include the following columns:
	
	\begin{tabularx}{\textwidth}{l l X}
		\toprule
		\textbf{Column} & \textbf{Required} & \textbf{Notes} \\
		\midrule
		% Fill in the actual columns once confirmed from source code
		name   & Yes & Full name of the client \\
		phone  & No  & Contact number \\
		email  & No  & Contact email \\
		\bottomrule
	\end{tabularx}
	
	\begin{steps}
		\item Navigate to the \ui{Import} section.
		\item Click \ui{Choose File} and select your CSV file.
		\item Review the preview.
		\item Click \ui{Import}.
	\end{steps}
	
	\dangerbox{Records created by import cannot be automatically undone.
		Make an export backup before importing new data.}
	
	\section{Generating PDF Documents}
	% How to download a service order as a PDF.
	% What the PDF contains (fields, layout).
	% Where the file is saved.
	
	\begin{steps}
		\item Open a service order.
		\item Click \ui{Download}.
		\item Choose a save location.
	\end{steps}
	
	% =============================================================
	%  CHAPTER 9 --- TROUBLESHOOTING
	% =============================================================
	\chapter{Troubleshooting}
	
	\section{Common Issues}
	
	\begin{tabularx}{\textwidth}{l X}
		\toprule
		\textbf{Problem} & \textbf{Solution} \\
		\midrule
		Application does not open
		& Check that WebView2 is installed (Windows only).
		Try reinstalling the application. \\
		\addlinespace
		CSV import fails
		& Verify the file is UTF-8 encoded and that all required
		columns are present with the correct header names. \\
		\addlinespace
		Data appears missing after reinstall
		& The database is stored separately from the application.
		See Section~\ref{sec:dataLocation} for the file location. \\
		\addlinespace
		PDF does not download
		& Check that you have write permission to the chosen folder. \\
		\bottomrule
	\end{tabularx}
	
	\section{Data File Location}
	\label{sec:dataLocation}
	% Where the SQLite database file is stored on each platform.
	% Windows: %APPDATA%\mechanic_shop\
	% macOS:   ~/Library/Application Support/mechanic_shop/
	% Linux:   ~/.local/share/mechanic_shop/
	
	\notebox{You can back up your data by copying the database file to a safe
		location. The filename is \texttt{mechanic\_shop.db}.}
	
	% =============================================================
	%  APPENDICES
	% =============================================================
	\appendix
	
	\chapter{Glossary}
	
	\begin{tabularx}{\textwidth}{l X}
		\toprule
		\textbf{Term} & \textbf{Definition} \\
		\midrule
		Client        & A person or entity whose vehicle is serviced. \\
		Vehicle       & A car or other vehicle registered under a client. \\
		Service Order & A record of a service performed or requested. \\
		CSV           & Comma-Separated Values. A plain text file format
		for tabular data. \\
		PDF           & Portable Document Format. A fixed-layout document
		format used for service records. \\
		\bottomrule
	\end{tabularx}
	
	\chapter{Keyboard Shortcuts}
	
	\begin{tabularx}{\textwidth}{l X}
		\toprule
		\textbf{Shortcut} & \textbf{Action} \\
		\midrule
		% Fill in once confirmed from source code
		\bottomrule
	\end{tabularx}
	
\end{document}